\documentclass[12pt,a4paper]{article}
\usepackage{fontspec}
\usepackage{float}
\usepackage{polyglossia}
\setotherlanguage{english}
\usepackage{amsmath,amssymb}
\usepackage{graphicx}
\usepackage{hyperref}
\usepackage{enumitem}
\usepackage[table]{xcolor}
\usepackage{array}
\usepackage{tabularx}
\usepackage{fancyhdr}
\usepackage{lastpage} % (tuỳ chọn)
\usepackage{listings}
\usepackage{xcolor}

\lstdefinestyle{sqlstyle}{
    language=SQL,
    basicstyle=\normalfont\small,            % <- bỏ \ttfamily để về font bình thường
    frame=single,
    framerule=0.4pt,
    breaklines=true,
    columns=fullflexible,
    showstringspaces=false,
    tabsize=2,
    numbers=left,
    numberstyle=\tiny,
    stepnumber=1,
    numbersep=8pt,
    backgroundcolor=\color{gray!8},
    keywordstyle=\bfseries,
    commentstyle=\itshape
}


\renewcommand{\arraystretch}{1.5} % tăng chiều cao dòng
\setlength{\tabcolsep}{10pt}      % tăng padding ngang

\title{MySQL - HackerRank Notebook}
\author{Pham Dang Hoang Thien}
\date{\today}

\begin{document}
% Footer cho TRANG 1
\fancypagestyle{titlefooter}{
    \fancyhf{} % xoá header/footer
    \renewcommand{\headrulewidth}{0pt}
    \renewcommand{\footrulewidth}{0pt}
    \fancyfoot[C]{%
        Student, Ho Chi Minh City, Vietnam\\
        e-mail: hoangthien312006@gmail.com%
    }
}

\maketitle
\thispagestyle{titlefooter} % trang 1: hiện footer, không hiện số trang

\clearpage                 % sang trang 2
\pagenumbering{arabic}     % kiểu 1,2,3...
\setcounter{page}{1}       % trang 2 sẽ mang số 1
\tableofcontents
\newpage
\section{Introduction to MySQL}
\section{Basic SQL Queries}
\subsection{Triangle Classification Query}
Write a query identifying the type of each record in the \texttt{TRIANGLES} table using its three side lengths.  
Output one of the following statements for each record in the table:

\begin{itemize}
    \item \textbf{Equilateral}: It's a triangle with 3 sides of equal length.
    \item \textbf{Isosceles}: It's a triangle with 2 sides of equal length.
    \item \textbf{Scalene}: It's a triangle with 3 sides of differing lengths.
    \item \textbf{Not A Triangle}: The given values of A, B, and C don't form a triangle.
\end{itemize}
\textbf{Input Format:}
\begin{table}[H]
\renewcommand{\arraystretch}{1.3}
\begin{tabular}{|c|c|}
\hline
\textit{Column} & \textit{Type} \\ \hline
\textit{A} & \textit{Integer} \\ \hline
\textit{B} & \textit{Integer} \\ \hline
\textit{C} & \textit{Integer} \\ \hline
\end{tabular}
\end{table}
The \texttt{TRIANGLES} table contains three columns: A, B, and C, representing the lengths of the sides of a triangle.
\newline
\noindent
\begin{minipage}[t]{0.48\textwidth}
\textbf{Sample Input:}
\vspace{0.3em}
\begin{tabular}{|c|c|c|}
\hline
A & B & C \\ \hline
20 & 20 & 23 \\ \hline
20 & 20 & 20 \\ \hline
20 & 21 & 22 \\ \hline
13 & 14 & 30 \\ \hline
\end{tabular}
\end{minipage}
\hfill
\begin{minipage}[t]{0.48\textwidth}
\textbf{Sample Output:}
\vspace{0.3em}
\begin{verbatim}
Isosceles
Equilateral
Scalene
Not A Triangle
\end{verbatim}
\end{minipage}

\subsection*{Solution Hint}
\textbf{Main idea:}  
Use a \texttt{CASE} expression to classify each triangle in this order:
\begin{enumerate}
    \item Check \textbf{Not A Triangle} first using the triangle inequality:
    \[
    A+B \le C \;\; \text{or} \;\; A+C \le B \;\; \text{or} \;\; B+C \le A
    \]
    \item If valid, check \textbf{Equilateral}: \texttt{A = B AND B = C}
    \item Check \textbf{Isosceles}: exactly two sides are equal
    \item Otherwise, it is \textbf{Scalene}
\end{enumerate}

\textbf{Required syntax / clauses:}
\begin{itemize}
    \item \texttt{CASE WHEN ... THEN ... ELSE ... END}
    \item Logical operators: \texttt{AND}, \texttt{OR}
    \item Comparison operators: \texttt{=}, \texttt{<=}
\end{itemize}

\subsection*{SQL Query}
\begin{lstlisting}[style=sqlstyle]
SELECT
  CASE
    WHEN A + B <= C OR A + C <= B OR B + C <= A THEN 'Not A Triangle'
    WHEN A = B AND B = C THEN 'Equilateral'
    WHEN A = B OR B = C OR A = C THEN 'Isosceles'
    ELSE 'Scalene'
  END AS triangle_type
FROM TRIANGLES;
\end{lstlisting}

\subsection*{Why this condition order?}
\begin{itemize}
    \item Put \texttt{Not A Triangle} first to filter out invalid side combinations.
    \item Then classify only valid triangles as equilateral, isosceles, or scalene.
\end{itemize}
\subsection{Revising Aggregations}
\textbf{Input Format:}
The CITY table is described as follows:
\begin{table}[H]
\renewcommand{\arraystretch}{1.3}
\begin{tabular}{|c|c|}
\hline
\textit{Fields} & \textit{Type} \\ \hline
\textit{ID} & \textit{Number} \\ \hline
\textit{NAME} & \textit{VARCHAR2(17)} \\ \hline
\textit{COUNTRYCODE} & \textit{VARCHAR2(3)} \\ \hline
\textit{DISTRICT} & \textit{VARCHAR2(20)} \\ \hline
\textit{POPULATION} & \textit{Number} \\ \hline
\end{tabular}
\end{table}
\subsubsection{The Count Function}
Query a count of the number of cities in CITY having a Population larger than 100,000.
\subsubsection*{Solution Hint}
The query should count the number of rows in the CITY table where the POPULATION is greater than 100,000.
\newline
\textit{\textbf{COUNT(*)}} is used to count all rows that meet the specified condition in the WHERE clause.
AS \textit{\textbf{city\_count}} renames the output column for clarity.
\subsubsection*{SQL Query}
\begin{lstlisting}[style=sqlstyle]
SELECT COUNT(*) AS city_count
FROM CITY
WHERE POPULATION > 100000;
\end{lstlisting}
\subsubsection{The Sum Function}
Query the total population of all cities in CITY where District is California.
\subsection*{Solution Hint}
The query should sum the POPULATION column in the CITY table where the DISTRICT is 'California'.
\newline
\textit{\textbf{SUM(POPULATION)}} is used to calculate the total population of rows that meet the specified condition in the WHERE clause.
AS \textit{\textbf{total\_population}} renames the output column for clarity.
\subsubsection*{SQL Query}
\begin{lstlisting}[style=sqlstyle]
SELECT SUM(POPULATION) AS total_population
FROM CITY
WHERE DISTRICT = 'California';
\end{lstlisting}
\subsubsection{Averages}
Query the average population of all cities in CITY where District is California.
\subsection*{Solution Hint}
The query should calculate the average of the POPULATION column in the CITY table where the DISTRICT is 'California'.
\newline
\textit{\textbf{AVG(POPULATION)}} is used to calculate the average population of rows that meet the specified condition in the WHERE clause.
AS \textit{\textbf{average\_population}} renames the output column for
clarity.
\subsubsection*{SQL Query}
\begin{lstlisting}[style=sqlstyle]
SELECT AVG(POPULATION) AS average_population
FROM CITY
WHERE DISTRICT = 'California';
\end{lstlisting}
\subsubsection{Average Population}
Query the average population for all cities in CITY, rounded down to the nearest integer.
\subsection*{Solution Hint}
The query should calculate the average of the POPULATION column in the CITY table and round it down to the nearest integer.
\begin{itemize}
    \item Use \textit{\textbf{AVG(POPULATION)}} to get the average population.
    \item Use \textit{\textbf{ROUND(...)}} to round the average population to the nearest integer.
    \item AS \textit{\textbf{average\_population}} renames the output column for clarity.
\end{itemize}
\subsubsection*{SQL Query}
\begin{lstlisting}[style=sqlstyle]
SELECT ROUND(AVG(POPULATION)) AS average_population
FROM CITY;
\end{lstlisting}
\subsubsection{Japan Population}
Query the total population of cities in CITY where CountryCode is JPN.
\subsection*{Solution Hint}
The query should sum the POPULATION column in the CITY table where the COUNTRYCODE is 'JPN'.
\newline
\textit{\textbf{SUM(POPULATION)}} is used to calculate the total population
of rows that meet the specified condition in the WHERE clause.
AS \textit{\textbf{total\_population}} renames the output column for clarity.
\subsubsection*{SQL Query}
\begin{lstlisting}[style=sqlstyle]
SELECT SUM(POPULATION) AS total_population
FROM CITY
WHERE COUNTRYCODE = 'JPN';
\end{lstlisting}
\subsubsection{Population Density Difference}
Query the difference between the maximum and minimum populations in CITY.
\subsection*{Solution Hint}
The query should calculate the difference between the maximum and minimum values of the POPULATION column in the
CITY table.
\begin{itemize}
    \item Use \textit{\textbf{MAX(POPULATION)}} to get the maximum population.
    \item Use \textit{\textbf{MIN(POPULATION)}} to get the minimum population.
    \item AS \textit{\textbf{pop\_diff}} renames the output column for clarity.
\end{itemize}

\subsubsection*{SQL Query}
\begin{lstlisting}[style=sqlstyle]
SELECT MAX(POPULATION) - MIN(POPULATION) AS pop_diff
FROM CITY;
\end{lstlisting}
\subsection{The Blunder}
Samantha was tasked with calculating the average monthly salaries for all employees in the EMPLOYEES table, but did not realize her keyboard's 0 key was broken until after completing the calculation. She wants your help finding the difference between her miscalculation (using salaries with any zeros removed), and the actual average salary.
Write a query calculating the amount of error (i.e.: \textit{actual-miscalculated} average monthly salaries), and round it up to the next integer.
\subsection*{Input Format}

The \texttt{EMPLOYEES} table is described as follows:

\begin{table}[H]
\centering
\renewcommand{\arraystretch}{1.3}
\begin{tabular}{|c|c|c|}
\hline
\textit{Column} & \textit{Type} \\ \hline
\textit{ID} & \textit{Integer} \\ \hline
\textit{Name} & \textit{String} \\ \hline
\textit{Salary} & \textit{Integer} \\ \hline
\end{tabular}
\end{table}

\textbf{Note:} Salary is per month.

\subsection*{Constraints}
\[
1000 < \text{Salary} < 10^5
\]

\subsection*{Sample Input}

\begin{table}[H]
\centering
\renewcommand{\arraystretch}{1.3}
\begin{tabular}{|c|c|c|}
\hline
\textit{ID} & \textit{Name} & \textit{Salary} \\ \hline
1 & Kristeen & 1420 \\ \hline
2 & Ashley   & 2006 \\ \hline
3 & Julia    & 2210 \\ \hline
4 & Maria    & 3000 \\ \hline
\end{tabular}
\end{table}

\subsection*{Sample Output}
\begin{verbatim}
2061
\end{verbatim}

\subsection*{Explanation}

The table below shows the salaries without zeros as they were entered by Samantha:

\begin{table}[H]
\centering
\renewcommand{\arraystretch}{1.3}
\begin{tabular}{|c|c|c|}
\hline
\textit{ID} & \textit{Name} & \textit{Salary} \\ \hline
1 & Kristeen & 142 \\ \hline
2 & Ashley   & 26  \\ \hline
3 & Julia    & 221 \\ \hline
4 & Maria    & 3   \\ \hline
\end{tabular}
\end{table}

Samantha computes an average salary of \(98.00\).  
The actual average salary is \(2159.00\).

The resulting error between the two calculations is
\[
2159.00 - 98.00 = 2061.00
\]
Since it is equal to the integer \(2061\), it does not get rounded up.
\subsection*{Solution Hint}
The query should calculate the difference between the actual average salary and the miscalculated average salary (with zeros removed). The difference should be rounded up to the next integer.
\begin{itemize}
    \item \textit{\textbf{AVG(Salary)}} is used to calculate the actual average salary.
    \item \textit{\textbf{REPLACE(CAST(Salary AS CHAR), '0', '')}} is used to remove all zeros from the salary by converting it to a string, replacing '0' with an empty string, and then casting it back to an unsigned integer.
    /dễ hiểu là bỏ hết số 0 trong salary (mô phỏng bàn phím hỏng) đi rồi tính trung bình.
    \item \textit{\textbf{AS UNSIGNED}} ensures that the modified salary is a positive integer for the average calculation. /nhận giá trị từ 0 trở lên.
\end{itemize}
\begin{lstlisting}[style=sqlstyle]
SELECT CEIL(
        AVG(Salary) - AVG(CAST(REPLACE(CAST(Salary AS CHAR), '0', '') AS UNSIGNED))
    ) AS error
FROM EMPLOYEES;
\end{lstlisting}

\end{document}